% =============================================================================
% paper_orbitalif.tex
%
% ORBITALIF: Orbital Attention-enhanced Leaky Integrate-and-Fire framework
%            for decentralized onboard federated optical cloud removal.
%
% NEW DRAFT (user-supplied abstract + introduction).  This file is intended
% to be the starting point for the user's own further editing; method,
% experiments, and conclusion sections are NOT included here yet.
%
% --- Reference verification convention --------------------------------------
%   Every \bibitem below is tagged in a comment with one of:
%     [VERIFIED]      — author / title / venue / year cross-checked against
%                       the actual primary source; OK to cite as-is.
%     [USER VERIFY]   — paper exists but exact bibliographic fields (page
%                       numbers, volume, doi) need user cross-check.
%     [PLACEHOLDER]   — generic stub; user must supply a real reference
%                       before submission.
% --- IMPORTANT --------------------------------------------------------------
%   Do NOT submit this paper while any \bibitem is still tagged [PLACEHOLDER]
%   or while [USER VERIFY] entries have not been confirmed against the
%   published or arXiv version.
% =============================================================================

\documentclass[conference]{IEEEtran}

\usepackage{cite}
\usepackage{amsmath,amssymb,amsfonts}
\usepackage{algorithmic}
\usepackage{graphicx}
\usepackage{textcomp}
\usepackage{xcolor}
\usepackage{hyperref}
\usepackage{booktabs}
\usepackage{caption}
\hypersetup{
  colorlinks   = true,
  linkcolor    = blue,
  citecolor    = blue,
  urlcolor     = blue,
}

% --- Useful macros ---
\newcommand{\E}{\mathbb{E}}
\newcommand{\R}{\mathbb{R}}
\newcommand{\D}{\mathcal{D}}
\newcommand{\Dir}{\mathrm{Dir}}
\newcommand{\OrbitALIF}{\textsc{OrbitALIF}}

\graphicspath{{figures/}}

\begin{document}

% --- Title block ---
\title{\OrbitALIF: Orbital Attention-enhanced Leaky Integrate-and-Fire
       for Decentralized Onboard Federated Optical Cloud Removal}

\author{
  \IEEEauthorblockN{(authors anonymised for review)}\\
  \IEEEauthorblockA{(institutions anonymised for review)}
}

\maketitle

% =====================================================================
% Abstract
% =====================================================================
\begin{abstract}
Satellite networks can acquire massive space information with advanced
remote sensing technologies, providing critical support for real-time
applications such as natural disaster monitoring. Emerging Space
Computing Power Networks provide a new architecture for coordinating
onboard computing resources, enabling satellites to perform intelligent
in-orbit processing collaboratively. Nevertheless, optical cloud
removal is a high-fidelity dense restoration task, and existing
artificial-neural-network-based models usually incur high computation
and energy costs, making them difficult to sustain under the limited
power supply and thermal conditions of satellite platforms.
To tackle this issue, we propose \OrbitALIF{}, which exploits the
sparse event-driven computation of spiking neural networks to reduce
inference energy, while enhancing pixel-level restoration capability
through attention- and receptive-field-enhanced modules and
architecture.
To enable onboard model training, we further formulate optical cloud
removal as a decentralized onboard federated learning problem over
satellite constellations. To the best of our knowledge, this is the
first work to study this problem. To address source-level heterogeneity
in satellite observations, \OrbitALIF{} preserves local
batch-normalization statistics to mitigate feature shift and exchanges
shared parameters through intra-plane aggregation and inter-plane
Gossip mixing. Experiments show that the $2.30$\,M onboard deployment
variant improves cloud-removal quality while reducing per-image
inference energy by up to $72.3\!\times$. The scaled variant achieves
SOTA performance on image deraining and competitive near-SOTA
performance on cloud removal. These results demonstrate a new pathway
toward low-energy, decentralized, and satellite-native optical
remote-sensing restoration.
\end{abstract}

\begin{IEEEkeywords}
Spiking neural networks, optical cloud removal, decentralized
federated learning, onboard satellite computing.
\end{IEEEkeywords}

% =====================================================================
% Section I: Introduction
% =====================================================================
\section{Introduction}\label{sec:intro}

Low-Earth-orbit optical remote-sensing satellite networks
continuously acquire large-scale space information, providing critical
data support for real-time applications such as natural disaster
monitoring, emergency response, agricultural observation, environmental
assessment, and defense surveillance~\cite{ref:earthobs1,ref:earthobs2}.
However, optical remote-sensing imagery is highly vulnerable to cloud
contamination, which degrades image usability and affects downstream
tasks such as object detection, land-cover classification, and change
detection. Conventional processing pipelines usually follow a
``satellite acquisition--raw data downlink--ground-server processing''
paradigm, where cloud-contaminated images are first transmitted to
ground stations and then restored and analyzed on ground servers. As
satellite constellations scale up and imaging resolution increases,
this paradigm faces an increasingly severe timeliness bottleneck:
visible contact windows between satellites and ground stations are
limited, downlink bandwidth is constrained, and ground-side processing
further introduces additional latency. Existing studies on onboard
federated learning have also pointed out that satellite data
transmission is a critical bottleneck for Earth-observation
applications, directly limiting real-time space information
services~\cite{ref:onboardfl-bottleneck}.

Space Computing Power Networks (Space-CPN) provide a promising
architecture to alleviate the above bottleneck. Instead of treating
satellites merely as data acquisition and forwarding nodes, Space-CPN
emphasizes the coordinated scheduling of communication, computation,
and intelligent processing resources in space networks, enabling
satellites to undertake part of the in-orbit data processing
workload~\cite{ref:spacecpn1}. For optical cloud removal, this paradigm
is particularly meaningful. Cloud contamination is a fundamental
problem that must be addressed before optical remote-sensing data can
be reliably used by downstream tasks. If cloud removal can be performed
onboard, the downlink of low-quality raw observations can be reduced,
the response speed of time-critical tasks can be improved, and cleaner
inputs can be provided for subsequent interpretation. Recent studies on
Space-CPN and in-orbit computing also regard distributed onboard
computation as an important direction for future space information
processing~\cite{ref:spacecpn1}.

Nevertheless, onboard optical cloud removal is not a simple transfer
of ground-side models to satellites. First, optical cloud removal is a
high-fidelity dense image restoration task. It requires pixel-level
recovery of cloud boundaries, thin-cloud textures, ground structures,
and high-frequency details. Existing artificial neural network
(ANN)-based cloud-removal models typically rely on continuous-valued
activations and dense multiply-accumulate operations, resulting in
considerable computation and energy costs. In particular, recent
diffusion-based cloud-removal methods have achieved strong performance
on high-resolution datasets such as CUHK-CR through progressive texture
recovery, but their iterative inference process is inherently in
tension with the limited power supply, thermal dissipation, and
computing budget of onboard satellite platforms~\cite{ref:cuhkcr}.
Therefore, onboard cloud-removal models should not only pursue offline
PSNR/SSIM performance, but also consider inference energy, computation
cost, and deployment feasibility.

Spiking neural networks (SNNs) provide a computational paradigm
different from conventional ANNs for low-energy onboard intelligence.
SNNs transmit information through sparse event-driven spike activities,
allowing silent neurons to skip unnecessary computation and enabling
low-cost spike operations on neuromorphic hardware~\cite{ref:maass1997,
ref:loihi}. For LEO satellite platforms constrained by solar power,
payload energy consumption, and thermal control, such sparse
computation is highly attractive. Previous studies have explored the
potential of SNNs in low-level vision restoration tasks such as image
deraining~\cite{ref:esdnet}, and SNNs have also been applied to
decentralized federated classification in satellite
scenarios~\cite{ref:flsnn}. However, optical cloud removal imposes
stronger requirements on SNNs. Binary spikes can reduce energy
consumption but may lose the fine-grained activation representation
needed for pixel-level regression. Meanwhile, cloud occlusion is
spatially non-uniform and spectrally complex, requiring the model to
capture cloud boundaries, semi-transparent thin-cloud structures, and
high-frequency ground textures simultaneously. Thus, an onboard SNN
cloud-removal model must balance sparse low-energy computation with
high-fidelity restoration.

In addition to inference energy, onboard model training is also
challenging. Due to differences in orbital positions, observation
regions, sensor states, atmospheric conditions, and cloud types, local
data on different satellites usually exhibit significant source-level
heterogeneity. For example, some satellites may observe more thin-cloud
samples, whereas others may mainly encounter thick clouds or complex
cloud distributions. Centralizing all raw observations on the ground
for training would incur high communication overhead and also
contradict the trend toward autonomous in-orbit processing emphasized
by Space-CPN. Federated
learning~\cite{ref:fedavg,ref:fedbn} provides a natural
solution: each satellite trains a model on its local
data and exchanges model parameters rather than raw observations.
However, standard centralized federated learning typically relies on a
ground parameter server, whereas LEO constellations are characterized
by intermittent connectivity and constrained intra-/inter-plane
communication topologies. Decentralized inter-satellite collaborative
training~\cite{ref:gossip-boyd} is therefore more suitable for
constellation scenarios. Moreover, source-level heterogeneous data
induce feature distribution shifts across satellites, and globally
shared batch-normalization statistics may destroy local feature
calibration. FedBN preserves local BN statistics to mitigate
feature-shift non-IID, which matches the source-level heterogeneous
satellite cloud-removal setting considered in this
work~\cite{ref:fedbn}.

The closest related work is FLSNN~\cite{ref:flsnn}, which demonstrates
the feasibility of decentralized federated SNN training on a
Walker-Star constellation for the EuroSAT classification task. However,
directly applying this paradigm to onboard optical cloud removal faces
three key gaps. The first is the \emph{task gap}: classification
produces discrete labels, whereas cloud removal is dense pixel
regression with different gradient structures, loss functions, and
restoration objectives. The second is the \emph{model gap}: satellite
cloud removal requires the reconstruction of multi-scale textures,
cloud boundaries, and spectral details, while small spiking
classification networks are insufficient for high-fidelity
reconstruction. The third is the \emph{data-heterogeneity gap}: non-IID
satellite observations in practical cloud-removal scenarios are driven
more by source factors such as sensors, cloud thickness, atmospheric
conditions, and observation regions than by class-label imbalance.
These gaps indicate that onboard federated cloud removal requires a
co-design across the task, model, and system levels.

To address these challenges, this paper studies a new problem:
\emph{decentralized onboard federated optical cloud removal over
satellite constellations}. To the best of our knowledge, this is the
first work to formulate optical cloud removal under this setting.
Different from existing ground-centered cloud-removal methods, this
work focuses on remote-sensing restoration under onboard communication,
energy, and computation constraints. Different from existing onboard
federated learning studies, this work addresses high-fidelity dense
image restoration rather than relatively lightweight tasks such as
classification or compression. The goal is not merely to improve
offline ground-side restoration metrics, but to establish a unified
design framework that jointly considers cloud-removal quality,
inference energy, onboard computation, and decentralized training.

To this end, we propose \textbf{\OrbitALIF{}}, an
\emph{Orbital Attention-enhanced Leaky Integrate-and-Fire} framework
for onboard optical cloud removal. \OrbitALIF{} exploits the sparse
event-driven computation of SNNs to reduce inference energy, while
improving pixel-level restoration capability through attention- and
receptive-field-enhanced modules and architecture. At the model level,
\OrbitALIF{} introduces five-level quantized spikes to improve the
representational capability of conventional binary spikes, and combines
dual-frequency residual modeling, frequency-spatial attention, adaptive
gated fusion, and spectro-temporal spike enhancement to recover cloud
boundaries, thin-cloud textures, and high-frequency ground details. At
the system level, \OrbitALIF{} constructs a decentralized constellation
training pipeline, retaining local BN statistics~\cite{ref:fedbn} to
mitigate source-level feature shift and exchanging shared parameters
through intra-plane aggregation and inter-plane Gossip
mixing~\cite{ref:gossip-boyd}. FedBN and Gossip are not proposed here
as new federated optimization algorithms; rather, they are adopted as
system-level training strategies adapted to constellation communication
topology and source-level heterogeneous data.

The main contributions of this work are summarized as follows.

\begin{itemize}
\item \textbf{We formulate decentralized onboard federated optical
  cloud removal.} To the best of our knowledge, this work is the first
  to formulate optical cloud removal as a decentralized onboard
  federated learning problem over satellite constellations, unifying
  optical remote-sensing restoration, onboard intelligent processing,
  and inter-satellite collaborative learning in a single framework.

\item \textbf{We propose \OrbitALIF{} for low-energy onboard cloud
  removal.} \OrbitALIF{} combines attention-enhanced LIF dynamics with
  sparse spiking computation. Through five-level quantized spikes,
  dual-frequency residual modeling, frequency-spatial attention,
  adaptive gated fusion, and spectro-temporal spike enhancement,
  \OrbitALIF{} preserves pixel-level restoration capability under
  low-power computation.

\item \textbf{We design a decentralized constellation training pipeline
  for source-level heterogeneous data.} To handle distribution
  differences across satellite observations, \OrbitALIF{} retains local
  BN statistics to mitigate feature shift and exchanges shared
  parameters through intra-plane aggregation and inter-plane Gossip
  mixing, avoiding both large-scale raw-data downlink and dependence on
  a centralized ground parameter server.

\item \textbf{We validate the quality--energy advantage and cross-task
  restoration potential of \OrbitALIF{}.} Experiments show that the
  $2.30$\,M onboard deployment variant improves cloud-removal quality
  while reducing per-image inference energy by up to $72.3\!\times$.
  The $14$\,M scaled variant reaches SOTA performance on image
  deraining and achieves competitive near-SOTA performance on optical
  cloud removal, indicating that attention-enhanced LIF dynamics have
  broader potential for low-level vision restoration.
\end{itemize}

The remainder of this paper is organized as follows.
Section~\ref{sec:sysmodel} introduces the constellation model, the
federated learning formulation, and the source-level Dirichlet non-IID
partition. Section~\ref{sec:framework} presents the \OrbitALIF{}
architecture and decentralized training pipeline.
Section~\ref{sec:experiments} reports centralized cloud-removal
results, component ablations, federated training results, and energy
analysis. Section~\ref{sec:conclusion} concludes this paper.

% =====================================================================
% Section II / III / IV / V — TO BE COMPOSED BY THE USER
% =====================================================================
% =====================================================================
% Section II: System Model
% =====================================================================
\section{System Model}\label{sec:sysmodel}

\subsection{LEO Constellation and Inter-Satellite Links}\label{sec:sysmodel-constellation}

We consider a Walker-Star low-Earth-orbit (LEO) satellite constellation
denoted by
\begin{equation}
\mathcal{S}
\;=\;
\left\{ s_{i,k} : i \in \mathcal{N}, \; k \in \mathcal{K}_i \right\},
\label{eq:constellation}
\end{equation}
where $\mathcal{N}$ is the set of orbital planes, $N\!=\!|\mathcal{N}|$
is the number of planes, and $\mathcal{K}_i$ is the set of satellites
in the $i$-th plane.  For notational simplicity we assume each plane
contains the same number of satellites,
$|\mathcal{K}_i|\!=\!K$ for all $i\!\in\!\mathcal{N}$.
Each satellite $s_{i,k}$ is equipped with inter-satellite laser links
(ISLs) for in-constellation model exchange. According to the
connection pattern, ISLs are divided into two categories: intra-plane
links and inter-plane links.

Intra-plane links connect adjacent satellites within the same orbital
plane. Since satellites in the same plane have relatively stable
geometry, these links are quasi-stationary across orbital periods and
are suitable for reliable ring-style aggregation. Inter-plane links
connect satellites from different orbital planes. Their availability
is more time-varying because of relative orbital motion, visibility
constraints, and Doppler effects. Following the Walker-Star setting
adopted in prior decentralized satellite
learning~\cite{ref:flsnn}, we instantiate $N\!=\!5$ and $K\!=\!10$,
yielding $NK\!=\!50$ satellites in total. For inter-plane
communication, we adopt a fixed chain topology as a simplified mixing
graph.

\subsection{Federated Optical Cloud-Removal Problem}\label{sec:sysmodel-problem}

Each satellite $s_{i,k}$ holds a local paired image dataset
\begin{equation}
\mathcal{D}_{i,k}
\;=\;
\left\{ (c_n, y_n) \right\}_{n=1}^{|\mathcal{D}_{i,k}|},
\label{eq:dataset}
\end{equation}
where $c_n\!\in\!\mathbb{R}^{3 \times H \times W}$ denotes a cloudy
image tile and $y_n\!\in\!\mathbb{R}^{3 \times H \times W}$ denotes
the corresponding cloud-free reference tile.  Let
$g_x : \mathbb{R}^{3 \times H \times W} \!\rightarrow\!
\mathbb{R}^{3 \times H \times W}$ be the \OrbitALIF{} cloud-removal
model parameterized by $x \!\in\! \mathbb{R}^d$, where
$d \!\approx\! 2.30 \!\times\! 10^6$ for the onboard deployment
variant. Let $\ell(g_x(c), y)$ denote the pixel-level reconstruction
loss, defined in Section~\ref{sec:framework-loss}. The global training
objective over the whole constellation is to minimise the empirical
risk
\begin{equation}
\min_{x \in \mathbb{R}^d}
f(x)
\;=\;
\frac{1}{NK}
\sum_{i \in \mathcal{N}}
\sum_{k \in \mathcal{K}_i}
f_{i,k}(x),
\label{eq:globalobj}
\end{equation}
where the local empirical risk on satellite $s_{i,k}$ is
\begin{equation}
f_{i,k}(x)
\;=\;
\frac{1}{|\mathcal{D}_{i,k}|}
\sum_{(c, y) \in \mathcal{D}_{i,k}}
\ell\!\left( g_x(c), \, y \right).
\label{eq:localobj}
\end{equation}

A centralised solution would require all paired cloudy/clear image
tiles to be downlinked to ground servers for training. This
contradicts the bandwidth and latency constraints discussed in
Section~\ref{sec:intro}. We therefore adopt a fully onboard
decentralised training strategy, where satellites compute local
updates on their own data and exchange model parameters through ISLs
without uploading raw observations.

\subsection{Decentralised Onboard Training Procedure}\label{sec:sysmodel-procedure}

The training proceeds in $R$ global communication rounds. Round $0$
starts from a common random initialisation $x^{(0)}\!\in\!\mathbb{R}^d$
that is broadcast to every satellite once at deployment time. Each
subsequent round $r\!\in\!\{0, \ldots, R-1\}$ then consists of four
steps. The symbol $r$ here counts inter-satellite \emph{communication
rounds} and is distinct from the spike-train time horizon $T\!=\!4$
of the underlying SNN, which is internal to a single forward pass and
discussed in Section~\ref{sec:framework}.

\subsubsection*{S1: Local update}
At the beginning of round $r$, each satellite initialises its model
from its plane-level broadcast model:
\begin{equation*}
x_{i,k}^{(r, 0)} \;=\; x_i^{(r)},
\end{equation*}
with $x_i^{(0)}\!=\!x^{(0)}$ at round $r\!=\!0$.
Satellite $s_{i,k}$ then performs $E$ local AdamW steps on its dataset
$\mathcal{D}_{i,k}$:
\begin{equation}
x_{i,k}^{(r, e+1)}
\;=\;
\mathrm{AdamW}\!\left(
x_{i,k}^{(r, e)},
\,
\nabla \tilde{f}_{i,k}\!\left( x_{i,k}^{(r, e)} \right),
\,
\eta_r
\right),
\quad e\!=\!0, \ldots, E\!-\!1,
\label{eq:localstep}
\end{equation}
where $\nabla\tilde{f}_{i,k}$ is the mini-batch stochastic estimate
of $\nabla f_{i,k}$ and $\eta_r$ is the learning rate at round $r$.
The notation $\mathrm{AdamW}(\cdot,\cdot,\cdot)$ stands for the
abstract local-update operator implemented by the AdamW optimiser;
replacing it by SGD or any other optimiser leaves the rest of the
training procedure unchanged.

\subsubsection*{S2: Intra-plane aggregation}
After local updates, the $K$ satellites in each orbital plane perform
ring AllReduce over quasi-stationary intra-plane ISLs. This produces
the plane-level consensus model
\begin{equation}
x_i^{(r + \frac{1}{2})}
\;=\;
\frac{1}{K}
\sum_{k \in \mathcal{K}_i}
x_{i,k}^{(r, E)}.
\label{eq:intraplane}
\end{equation}
To accommodate source-level feature heterogeneity, we adopt the
\textsc{FedBN}~\cite{ref:fedbn} normalisation protocol: every
batch-normalisation parameter, including the affine
$(\gamma,\beta)$, the running mean, and the running variance, is
\emph{excluded} from the averaging in (\ref{eq:intraplane}); these
parameters are taken from the plane leader of plane $i$ verbatim and
are not averaged across the $K$ satellites of the plane. The same
exclusion is applied later in the inter-plane Gossip step
(Section~\ref{sec:sysmodel-procedure}, S3). Local training in S1
then refines the BN parameters of each satellite on its own data, so
that BN statistics naturally adapt to the source distribution
observed locally.

\subsubsection*{S3: Inter-plane Gossip mixing}
After intra-plane aggregation, each orbital plane obtains a plane-level
model $x_i^{(r+\frac{1}{2})}$. The $N$ plane-level models are then
mixed through single-step Gossip averaging~\cite{ref:gossip-boyd} over
the inter-plane chain topology:
\begin{equation}
x_i^{(r+1)}
\;=\;
\sum_{j \in \mathcal{N}}
W_{ij} \, x_j^{(r + \frac{1}{2})},
\quad i \in \mathcal{N},
\label{eq:gossip}
\end{equation}
where $W \!\in\! \mathbb{R}^{N \times N}$ is a symmetric doubly
stochastic mixing matrix satisfying
\begin{equation*}
W \mathbf{1} \!=\! \mathbf{1}, \quad
\mathbf{1}^{\!\top} W \!=\! \mathbf{1}^{\!\top}, \quad
W \!=\! W^{\!\top}.
\end{equation*}
Consistent with the \textsc{FedBN} protocol adopted in S2, all
batch-normalisation parameters (affine $(\gamma,\beta)$, running mean,
and running variance) are also excluded from the Gossip mixing in
(\ref{eq:gossip}); only the convolutional weights and other dense
layer weights are mixed across orbital planes. In this way, the
constellation shares restoration knowledge across orbital planes while
avoiding the destructive averaging of feature statistics that would
arise under pure FedAvg-style aggregation~\cite{ref:fedavg} when
satellites observe different cloud thicknesses, sensors, and
atmospheric conditions.

\subsubsection*{S4: Intra-plane broadcast}
Finally, each orbital plane broadcasts the mixed model
$x_i^{(r+1)}$ to all satellites $s_{i,k}\!\in\!\mathcal{K}_i$ in the
same plane. This model is used as the initialisation for the next
communication round. The above four steps complete one round of
decentralised onboard training.

\subsection{Source-Level Dirichlet Non-IID Partition}\label{sec:sysmodel-dirichlet}

To simulate the natural source-level feature heterogeneity of
satellite imagery, we do not use the conventional label-level Dirichlet
partition commonly adopted in classification tasks. Instead, we
construct a source-level non-IID partition over the training pool.
Specifically, the joint training set is divided into $S\!=\!2$ source
domains: CUHK-CR1 represents thin-cloud scenes, and CUHK-CR2
represents thick-cloud scenes~\cite{ref:cuhkcr}. For each satellite
$s_{i,k}$, we draw a source mixing vector
\begin{equation}
\mathbf{p}_{i,k}
\;=\;
\left( p_{i,k}^{(1)}, \, p_{i,k}^{(2)} \right)
\;\sim\;
\Dir\!\left( \alpha \mathbf{1}_S \right).
\label{eq:dirichlet}
\end{equation}
Satellite $s_{i,k}$ then samples image tiles from the $s$-th source
domain in proportion to $p_{i,k}^{(s)}$. To ensure executable local
training, each satellite is assigned at least five training samples.

The Dirichlet concentration parameter $\alpha$ controls the strength
of source-level heterogeneity. A smaller $\alpha$ produces more
concentrated single-source clients, where some satellites mainly hold
thin-cloud samples while others mainly hold thick-cloud samples. A
larger $\alpha$ yields a more balanced source mixture and approaches
an IID-like setting. We use $\alpha\!=\!0.1$ as the default
configuration. Under this setting, $30$ out of $50$ satellites hold
only the minimum of five samples ($60\%$ of the constellation), while
three satellites individually hold more than $100$ samples, with the
actual partition log recording the three largest client sizes as
$\{117, 118, 118\}$. This partition reflects both the data sparsity
and source-level heterogeneity that may arise in onboard cloud-removal
scenarios. Sensitivity to $\alpha$ is reported in
Section~\ref{sec:exp-federated}.

% =====================================================================
% Section III: The ORBITALIF Framework
% =====================================================================
\section{The \OrbitALIF{} Framework}\label{sec:framework}

The \OrbitALIF{} framework consists of two co-designed components:
an attention-enhanced LIF restoration backbone
(Section~\ref{sec:framework-arch}) and a decentralized constellation
training pipeline (Section~\ref{sec:framework-fl}). Figure~\ref{fig:overview}
illustrates the overall system. On each satellite, the restoration
backbone takes a cloudy RGB image tile
$\mathbf{c}\in\mathbb{R}^{B\times 3\times H\times W}$ as input and
predicts a cloud-free image $\hat{\mathbf{y}}$ through a sequence of
spiking restoration modules. Across the constellation, satellites
collaboratively train the shared restoration parameters through
intra-plane aggregation and inter-plane Gossip mixing, while all
BN-tagged parameters are kept local following the \textsc{FedBN}
protocol described in Section~\ref{sec:sysmodel-procedure}.

\begin{figure*}[!htbp]
  \centering
  \includegraphics[width=\textwidth]{fig1.pdf}
  \caption{\OrbitALIF{} system overview. \textbf{(a)} Single-satellite
  restoration model: a three-level multi-scale spiking restoration
  backbone with full-resolution spectro-temporal spike enhancement,
  deeper dual-frequency residual modeling, adaptive gated cross-stage
  feature fusion, and frequency-spatial attention. The backbone predicts
  a residual cloud-removal image that is added to the cloudy input.
  \textbf{(b)} Decentralized $5\times10$ Walker-Star constellation:
  satellites perform intra-plane aggregation over quasi-stationary
  intra-plane ISLs and inter-plane Gossip mixing over the chain topology.
  All BN-tagged parameters remain local and are excluded from both
  aggregation stages.}
  \label{fig:overview}
\end{figure*}

\subsection{\OrbitALIF{} Restoration Backbone}
\label{sec:framework-arch}

\OrbitALIF{} adopts a three-level multi-scale spiking restoration
backbone designed for onboard optical cloud removal. The shallow level
preserves the full $H\times W$ spatial resolution to retain cloud
boundaries, thin-cloud textures, and high-frequency ground details.
The two deeper levels progressively reduce the spatial resolution by
PixelUnshuffle while increasing the feature width, allowing the model
to capture larger-context cloud structures under compact computation.
A symmetric synthesis path reconstructs the clean image from the
multi-scale spike features, and adaptive gated cross-stage fusion is
used to combine shallow spatial details with deeper restoration
responses.

Let $T$ denote the spike-train time horizon inside a single forward
pass. In this work, all spiking modules use $T=4$, which is distinct
from the federated communication-round count $R$ in
Section~\ref{sec:sysmodel-procedure}. Given a cloudy input image
$\mathbf{c}$, the backbone first maps it into multi-scale spike
features through the analysis path. The synthesis path then produces
a residual restoration image, which is added to the input through a
global residual connection:
\begin{equation}
\hat{\mathbf{y}}
=
\mathbf{c}
+
\mathrm{Conv}
\left(
\overline{
\mathrm{Synth}
\left(
\mathrm{Anal}(\mathbf{c})
\right)
}^{(T)}
\right),
\label{eq:globalresidual}
\end{equation}
where $\mathrm{Anal}(\cdot)$ and $\mathrm{Synth}(\cdot)$ denote the
analysis and synthesis paths, respectively, and
$\overline{(\cdot)}^{(T)}$ denotes temporal averaging over the
$T=4$ spike steps.

The deployment variant sets the base channel width to $C=24$,
resulting in approximately $2.30$\,M trainable parameters. This compact
configuration is used for centralized cloud-removal experiments,
ablation studies, energy analysis, and decentralized federated
training. We also instantiate a capacity-scaled variant with
$C=56$ and approximately $14$\,M parameters. This larger variant shares
the same dataflow and is used to study restoration scalability and the
quality--computation trade-off, rather than to claim an unconditional
cloud-removal SOTA result.

The backbone is constructed from five main components. The five-level
quantized spike (\FQS{}) activation improves the representation
granularity of sparse LIF dynamics. The dual-frequency residual block
(\DFRB{}) captures complementary temporal and spatial-frequency
responses through two cascaded processing groups, with frequency-spatial
attention applied at the tail of each block. The spectral-spatial
hybrid attention module (\SHAM{}) recalibrates spike features along
temporal and frequency-spatial dimensions and is embedded within each
\DFRB{}. The adaptive gated fusion module (\AGFM{}) performs
content-adaptive cross-stage feature fusion at the first and second
multi-scale levels. Finally, the spectro-temporal spike hyper-block
(\SSHB{}) enhances full-resolution features through temporal expansion,
spike-convolutional refinement, and frequency-domain spatial mixing.
These components are detailed in the following subsections.

\subsection{\FQS{}: Five-Level Quantized Spike}
\label{sec:framework-5qs}

Binary LIF spikes provide sparse event-driven computation, but they
restrict each time step to a one-bit activation state. This is overly
restrictive for pixel-level cloud removal, where intermediate features
must preserve weak color gradients, thin-cloud transitions, and
fine-grained texture variations. Replacing spikes with continuous
ANN-style activations would recover a richer dynamic range, but would
also weaken the sparse operation pattern that makes SNNs attractive
for onboard inference.

To balance activation expressiveness and event-driven sparsity,
\OrbitALIF{} introduces the five-level quantized spike (\FQS{})
activation. Instead of emitting a binary spike $o_t\in\{0,1\}$, a
\FQS{} neuron clamps the membrane potential $u_t$ to $[0,4]$, rounds
to the nearest integer, and normalizes by 4:
\begin{equation}
o_t
=
\frac{
\mathrm{round}\!\left(\mathrm{clamp}(u_t,\,0,\,4)\right)
}{4}
\;\in\;
\left\{0,\,\tfrac{1}{4},\,\tfrac{1}{2},\,\tfrac{3}{4},\,1\right\}.
\label{eq:5qs-forward}
\end{equation}
A synaptic operation is counted only when $o_t>0$, so the activation
remains event-driven while providing five discrete non-negative
amplitude levels above the silent state. In this sense, \FQS{} improves
per-step representation granularity without reverting to dense
continuous ANN activations.

The quantization in \eqref{eq:5qs-forward} is non-differentiable.
During training, we use a straight-through rectangular surrogate
gradient over the active quantization range $[0,4]$:
\begin{equation}
\frac{\partial o_t}{\partial u_t}
\approx
\frac{1}{4}\cdot
\mathbb{1}\bigl\{0\le u_t\le 4\bigr\}.
\label{eq:5qs-surrogate}
\end{equation}
This surrogate passes scaled gradients only within the effective
five-level activation window and stabilizes the optimization of
multi-level quantized spikes. In our ablation study, replacing binary
spikes with \FQS{} consistently improves cloud-removal fidelity,
indicating that multi-level spike amplitudes are important for
restoring subtle cloud and texture details.

\subsection{\DFRB{}: Dual-Frequency Residual Block}
\label{sec:framework-dfrb}

Optical cloud removal requires two complementary types of information:
temporal spike responses that describe cloud boundary transitions
across the $T$-step spike train, and spatial high-frequency responses
that preserve vegetation texture, urban grids, and fine-grained ground
details. A single local spiking convolution tends to emphasize only
one part of this spectrum. As a result, the model may recover broad
cloud contours while losing fine textures, or preserve local textures
while failing to model temporally diffuse cloud structures.

The dual-frequency residual block (\DFRB{}) is designed to capture
these complementary responses within a sparse spiking dataflow.
Rather than using two independent parallel branches, \DFRB{}
consists of two cascaded processing groups applied sequentially.
Figure~\ref{fig:overview}(c) schematically illustrates the dual-frequency
information flow; the two groups are applied sequentially in the
actual implementation.

The first group extracts temporal spike responses. Given the input
feature $\mathbf{X}$, a standard LIF neuron separates it into a
high-frequency temporal component $\mathbf{X}_{h}$ and a residual
low-frequency component $\mathbf{X}_{l}=\mathbf{X}-\mathbf{X}_{h}$.
These are combined with learnable scales $\alpha_h,\alpha_l$ and a
nonlinear interaction term before being passed through a convolutional
layer and BN:
\begin{equation}
\mathbf{Z}_1
=
\mathrm{BN}\!\left(
\mathrm{Conv}\!\left(
\alpha_h \mathbf{X}_{h}
+ \alpha_l \mathbf{X}_{l}
+ \mathbf{X}\odot\mathbf{X}_{h}
\right)\right).
\label{eq:dfrb-g1}
\end{equation}
The second group applies spatial high-frequency extraction to
$\mathbf{Z}_1$ using a PixelShuffle-LIF block
(\texttt{PixelShuffleLIFBlock}), which spatially rearranges the
feature map by $2\times$ pixel unshuffling and applies LIF neurons to
the rearranged spatial-temporal representation, yielding a spatial
high-frequency component $\mathbf{Z}_{2h}$ and a residual component
$\mathbf{Z}_{2l}=\mathbf{Z}_1-\mathbf{Z}_{2h}$.
A learnable scalar cross-scale gate $\sigma(g)$ combines the
high-frequency responses from both groups:
\begin{equation}
\mathbf{Z}_2
=
\mathrm{BN}\!\left(
\mathrm{Conv}\!\left(
\alpha_h' \mathbf{Z}_{2h}
+ \alpha_l' \mathbf{Z}_{2l}
+ \mathbf{Z}_1\odot\mathbf{Z}_{2h}
+ \sigma(g)\,\mathbf{X}_{h}\odot\mathbf{Z}_{2h}
\right)\right).
\label{eq:dfrb-g2}
\end{equation}
The gate $\sigma(g)$ is initialized to $g\!=\!0$, so it starts near
zero and gradually opens as the model learns when cross-group
high-frequency compensation is beneficial.

After the two cascaded groups, a shortcut connection is added and the
combined feature passes through a temporal-channel-spatial attention
module (TCS-Att) followed by the spectral-spatial hybrid attention
module (\SHAM{}, Section~\ref{sec:framework-sham}):
\begin{equation}
\mathrm{DFRB}(\mathbf{X})
=
\SHAM{}\!\left(
\mathrm{TCS\text{-}Att}\!\left(
\mathbf{Z}_2 + \mathrm{Conv}(\mathbf{X})
\right)
\right).
\label{eq:dfrb}
\end{equation}
The ablation study shows that removing the spatial-frequency component
of \DFRB{} leads to a clear drop in PSNR, confirming that explicit
dual-frequency modeling is important for thin-cloud and
complex-texture restoration.

\subsection{\SHAM{}: Spectral-Spatial Hybrid Attention}
\label{sec:framework-sham}

Spiking restoration features contain not only channel and spatial
dimensions, but also a temporal dimension induced by the spike-train
horizon. Conventional channel or spatial attention modules recalibrate
only one or two dimensions and do not explicitly exploit the temporal
responses of SNNs or the frequency structures caused by cloud
degradation. In cloud removal, thin clouds often introduce low-frequency
brightness shifts together with high-frequency texture attenuation,
while thick cloud boundaries create strong spatial discontinuities.
This motivates attention along both temporal-amplitude and
frequency-spatial axes.

Given a spike feature volume
$\mathbf{X}\in\mathbb{R}^{T\times C\times H\times W}$, \SHAM{}
contains two complementary branches applied in series. The
temporal-amplitude attention (TAA) branch first pools $\mathbf{X}$
over spatial dimensions independently for each time step, mixes the
resulting per-step descriptors through a learned linear projection over
the temporal axis, and produces per-step scalar weights that are
broadcast back to modulate each time step:
\begin{equation}
\mathbf{X}'
=
\mathbf{X}
+
\mathbf{X}
\odot
\sigma\!\left(
\mathrm{FC}_{T\to T}\!\left(
\overline{\mathbf{X}}_{\mathrm{sp}}
\right)
\right),
\label{eq:sham-taa}
\end{equation}
where $\overline{\mathbf{X}}_{\mathrm{sp}}$ is the spatial-pooled
descriptor and the FC layer captures inter-timestep amplitude
correlations.

The frequency-spatial attention (SFA) branch computes the temporal
mean of $\mathbf{X}'$, applies a two-dimensional real FFT to obtain
the magnitude spectrum, and uses a $1\times1$ convolutional MLP to
selectively weight each spatial frequency component while preserving
the original phase:
\begin{equation}
\mathbf{M}
=
\left|
\mathcal{F}_{\mathrm{rFFT2}}\!\left(
\overline{\mathbf{X}'}^{(T)}
\right)
\right|,
\label{eq:sham-rfft}
\end{equation}
\begin{equation}
\hat{\mathbf{M}}
=
\mathbf{M}
\odot
\sigma\!\left(\mathrm{FreqMLP}(\mathbf{M})\right).
\label{eq:sham-freqmlp}
\end{equation}
The modulated spectrum is reconstructed via inverse FFT (iFFT) with
the original phase restored, and a $7\times7$ spatial convolution
generates the frequency-spatial attention mask
$\mathbf{w}_S\in\mathbb{R}^{1\times H\times W}$.

The full \SHAM{} output is computed as a learnable gated residual,
where the scale $\eta$ is initialized to zero so the module starts as
an identity mapping:
\begin{equation}
\mathrm{SHAM}(\mathbf{X})
=
\mathbf{X}
+
\sigma(\eta)
\left(
\mathbf{X}'
\odot
\mathbf{w}_S
-
\mathbf{X}
\right).
\label{eq:sham}
\end{equation}
This design allows \OrbitALIF{} to exploit both the temporal response
patterns of SNNs and the frequency-domain structures of cloud
contamination. \SHAM{} is embedded at the tail of each \DFRB{} block
(after TCS-Att) and thus recalibrates the dual-frequency residual
output before it propagates to the next block.

\subsection{\AGFM{}: Adaptive Gated Fusion Module}
\label{sec:framework-agfm}

Multi-scale restoration backbones need to combine shallow spatial
details with deeper restoration responses. Direct addition or
concatenation applies the same fusion rule to every spatial location,
which is suboptimal for cloud removal. In thin-cloud or clear regions,
shallow features often contain reliable ground textures and should be
preserved. In thick-cloud regions, however, shallow features are
themselves corrupted, and the synthesis-path restoration response
should dominate.

The adaptive gated fusion module (\AGFM{}) performs content-adaptive
cross-stage fusion. Let $\mathbf{d}$ denote the upsampled feature from
the synthesis path and $\mathbf{e}$ the encoder skip feature from the
analysis path. \AGFM{} concatenates the two features along the channel
dimension and generates a spatially varying gate through a
$1\times1$ convolution:
\begin{equation}
\boldsymbol{\gamma}
=
\sigma\!\left(
\mathrm{Conv}_{1\times1}
\!\left(
[\mathbf{d};\mathbf{e}]_C
\right)
\right),
\label{eq:agfm-gate}
\end{equation}
where $[\cdot;\cdot]_C$ denotes channel-wise concatenation and
$\boldsymbol{\gamma}\in\mathbb{R}^{C\times H\times W}$. The fused
feature is
\begin{equation}
\mathbf{f}
=
\boldsymbol{\gamma}\odot\mathbf{d}
+
(1-\boldsymbol{\gamma})\odot\mathbf{e}.
\label{eq:agfm}
\end{equation}
The gate allows the model to dynamically select between shallow details
and deeper restoration responses according to local cloud conditions.
In \OrbitALIF{}, \AGFM{} is inserted at the first and second
multi-scale levels only. At the deepest level, the encoder and decoder
communicate through a direct transition without a separate cross-stage
fusion module.

\subsection{\SSHB{}: Spectro-Temporal Spike Hyper-Block}
\label{sec:framework-sshb}

The full-resolution level contains the most detailed spatial
information and is therefore critical for restoring cloud boundaries,
thin-cloud textures, and high-frequency ground structures. However,
full-resolution processing is also the most expensive part of the
network. A local spiking convolution has limited receptive field, while
direct global attention would introduce excessive onboard computation.
The spectro-temporal spike hyper-block (\SSHB{}) is designed to enhance
full-resolution features by combining temporal expansion,
spike-convolutional refinement, and frequency-domain spatial mixing.

Given an input spike feature $\mathbf{X}\in\mathbb{R}^{T\times C\times H\times W}$,
\SSHB{} first applies \DFRB{} to extract dual-frequency residual
responses with embedded frequency-spatial attention. The resulting
feature is then processed by a $2\times$ PixelUnshuffle operation,
which rearranges the spatial grid into a richer temporal-channel
context by expanding the effective time horizon from $T$ to $T\cdot r^2$
(with $r=2$, this yields $16$ steps at half resolution). A LIF node
applies event-driven gating on the expanded representation, and a
temporal-channel attention module (TimeAttention via 3D pooling and
shared MLP) recalibrates inter-step importance. A 3D convolution then
compresses the $T\cdot r^2$-step representation back to $T$ steps.

After temporal compression, two spike-convolutional refinement blocks
(each: LIF$\to$Conv$\to$BN$\to$TCS-Att) further process the
reduced-resolution features. The refined feature is then upsampled
back to the original spatial resolution and combined with the skip
feature produced by the initial \DFRB{} (before PixelUnshuffle) through
direct addition. A final TCS-Att head jointly modulates temporal,
channel, and spatial responses. Finally, a frequency-spectral
enhancement block (FSE) applies a lightweight MLP in the 2D FFT
magnitude domain while preserving phase, providing global-receptive-field
feature mixing in $O(N\log N)$ without dense spatial attention.

The overall \SSHB{} dataflow is:
\begin{equation}
\mathrm{SSHB}(\mathbf{X})
=
\mathrm{FSE}
\circ
\mathrm{TCS\text{-}Att}
\circ
(\cdot + \mathrm{DFRB}(\mathbf{X}))
\circ
\mathrm{Refine}
\circ
\mathrm{TCAM}
\circ
\mathrm{PU}_{2}
\circ
\mathrm{DFRB}(\mathbf{X}),
\label{eq:sshb}
\end{equation}
where $\mathrm{PU}_{2}$ denotes $2\times$ PixelUnshuffle followed by
LIF and temporal compression, $\mathrm{TCAM}$ denotes the TimeAttention
module, $\mathrm{Refine}$ denotes the two spike-convolutional refinement
blocks, and $(\cdot + \mathrm{DFRB}(\mathbf{X}))$ denotes the skip
addition before the final TCS-Att head. \SSHB{} is placed at the
full-resolution level where pixel-level details are most important.
Although each \SSHB{} block is more expressive than a local residual
block, its cost is controlled by 5QS-LIF sparsity and FFT-based
spatial mixing.

\subsection{Training Objective}
\label{sec:framework-loss}

\OrbitALIF{} is trained end-to-end using a composite reconstruction
objective that combines a Charbonnier penalty with a structural
similarity term. Given the prediction $\hat{\mathbf{y}}$ and the
cloud-free reference image $\mathbf{y}$, the loss is defined as
\begin{equation}
\ell(\hat{\mathbf{y}},\mathbf{y})
=
\frac{1}{|\Omega|}
\sum_{p\in\Omega}
\sqrt{
(\hat{\mathbf{y}}_p-\mathbf{y}_p)^2+\varepsilon^2
}
+
\lambda
\left(
1-\mathrm{SSIM}(\hat{\mathbf{y}},\mathbf{y})
\right),
\label{eq:loss}
\end{equation}
where $\Omega$ is the pixel set, $\varepsilon=10^{-3}$, and
$\lambda=0.1$. The Charbonnier term acts as a smooth robust
reconstruction loss and reduces the influence of local outliers caused
by severe cloud occlusion. The SSIM term encourages structural
consistency, helping the model preserve cloud boundaries, textures, and
local image structures. With $\lambda=0.1$, the Charbonnier term
dominates the optimization throughout training.

\subsection{Decentralized Constellation Training Pipeline}
\label{sec:framework-fl}

Section~\ref{sec:sysmodel-procedure} defines the mathematical procedure
of decentralized onboard training. This subsection specifies the
implementation choices used to instantiate that procedure for
\OrbitALIF{}.

First, \OrbitALIF{} adopts a \textsc{FedBN}-style parameter handling
protocol to address source-level feature heterogeneity. Under the
source-level Dirichlet partition, different satellites may observe
different cloud-density regimes, which induces distinct feature
statistics and spike activation patterns. Averaging BN parameters
globally would force these heterogeneous distributions into a single
normalization state and may destabilize both intermediate features and
spike firing rates. Therefore, all BN-tagged parameters are kept local
and excluded from both intra-plane aggregation and inter-plane Gossip.
This includes BN affine parameters $(\gamma,\beta)$, running means,
running variances, batch counters, and the corresponding TDBN and
TemporalBN parameters used by spiking modules. Only the shared non-BN
restoration parameters — convolutional weights, biases, and
non-BN attention parameters — are exchanged across satellites. Local
training in step S1 of each round then naturally refines the BN
statistics of each satellite on its own data, so that BN normalization
adapts to the local source distribution without cross-satellite
interference.

Second, inter-plane communication is implemented by single-step Gossip
mixing over a fixed chain topology. Compared with a ground-server-based
parameter aggregation scheme, Gossip better matches the intermittent
connectivity and sparse inter-plane topology of LEO constellations.
Each orbital plane exchanges shared non-BN parameters only with its
neighbors according to the doubly stochastic mixing matrix $W$ defined
in Section~\ref{sec:sysmodel-procedure}. \textsc{FedBN}~\cite{ref:fedbn}
and Gossip averaging~\cite{ref:gossip-boyd} are existing methods; they
are not proposed here as new federated optimization algorithms. Rather,
they are adopted as system-level training strategies that are adapted
to the constellation communication topology and source-level
heterogeneous cloud-removal data.

Third, communication accounting is performed only over shared non-BN
parameters. For the $2.30$\,M onboard deployment variant, inter-plane
Gossip transfers the non-BN parameter subset, with each 32-bit
floating-point state transfer requiring approximately
$4\,d_{\mathrm{shared}}$ bytes, where $d_{\mathrm{shared}}$ is the
number of shared parameters. Since the inter-plane topology is a
fixed chain, the per-round cross-plane traffic scales with the number
of active inter-plane edges and the shared parameter count. The
concrete per-round and cumulative communication costs are reported in
Section~\ref{sec:exp-federated}.

Overall, the decentralized training pipeline has three properties: raw
cloudy observations remain local to satellites, BN-tagged parameters
preserve local feature calibration under source-level heterogeneity,
and shared restoration parameters are collaboratively updated through
intra-plane aggregation and inter-plane Gossip mixing. This allows
\OrbitALIF{} to train onboard without uploading raw observations or
relying on a centralized ground parameter server.

% =====================================================================
% Section IV: Experiments
% =====================================================================
\section{Experiments}\label{sec:experiments}

\subsection{Experimental Setup}\label{sec:exp-setup}

\textbf{Datasets.}\;\;We use the publicly available CUHK-CR
benchmark~\cite{ref:cuhkcr}, which contains paired cloudy and
cloud-free Sentinel-2 RGB tiles at high resolution. CUHK-CR1 contains
thin-cloud scenes and CUHK-CR2 contains thick-cloud scenes; the two
splits define the two source domains $\{1,2\}$ used by the
source-level Dirichlet partition in
Section~\ref{sec:sysmodel-dirichlet}. Following the original benchmark
protocol, all training and evaluation are performed on
$64\!\times\!64$ random crops, with PSNR and SSIM evaluated on the
held-out test split using $64\!\times\!64$ centre patches.

\textbf{Models.}\;\;Three matched-protocol backbones are evaluated.
\OrbitALIF{}-SNN is the spiking variant described in
Section~\ref{sec:framework-arch} with five-level quantized spikes,
TDBN normalization, and $T\!=\!4$ spike steps. \OrbitALIF{}-ANN keeps
the identical dataflow but replaces every \FQS{}-LIF activation with
ReLU and uses standard 2D BatchNorm; this isolates the effect of
spiking computation from the effect of architectural choices.
PlainUNet is a parameter-light vanilla U-Net baseline ($1.31$\,M
parameters) commonly used for cloud-removal ablation. Unless
otherwise stated, the spiking variant uses base channel width
$C\!=\!24$ ($\approx\!2.30$\,M parameters) and is referred to as the
\emph{onboard deployment variant}.

\textbf{Optimization.}\;\;All centralized models are trained for
$600$ epochs with AdamW, base learning rate $10^{-3}$, cosine decay
to $10^{-7}$, $3$-epoch linear warmup, batch size $4$, and gradient
clipping at $\|g\|_2\!\le\!1$. The composite reconstruction loss in
\eqref{eq:loss} uses $\varepsilon\!=\!10^{-3}$ and
$\lambda\!=\!0.1$. PlainUNet uses the same optimizer settings but
converges in $400$ epochs. Ablation runs (Section~\ref{sec:exp-ablation})
use $300$ epochs to keep training cost tractable.

\textbf{Federated training.}\;\;We instantiate the constellation of
Section~\ref{sec:sysmodel-constellation} with $N\!=\!5$ orbital planes
and $K\!=\!10$ satellites per plane ($NK\!=\!50$ in total). Federated
training proceeds for $R\!=\!200$ communication rounds. Within each
round, every satellite performs $E\!=\!2$ AdamW local epochs over its
Dirichlet-partitioned data shard, intra-plane ring AllReduce uses
$2$ averaging passes, and inter-plane Gossip uses single-step mixing
over a fixed chain mixing matrix~$W$. The default Dirichlet
concentration is $\alpha\!=\!0.1$.

\textbf{Hardware and metrics.}\;\;All training and inference experiments
use a single NVIDIA RTX-class GPU. We report PSNR (dB) and SSIM
averaged over the held-out CUHK-CR1 / CR2 test splits. For the
federated runs we additionally report cumulative inter-plane
communication volume in megabytes and total wall time in hours. For
energy analysis (Section~\ref{sec:exp-energy}) we use forward hooks to
measure per-layer MAC counts and non-zero activation rates on
$32$ test patches and combine these measurements with literature
per-operation energy constants.

\subsection{Centralized Cloud-Removal Performance}\label{sec:exp-centralized}

Table~\ref{tab:centralized} reports centralized results on CUHK-CR1
and CUHK-CR2. \OrbitALIF{} attains
$25.374$\,dB\,/\,$0.7728$ PSNR\,/\,SSIM on CR1 and
$23.197$\,dB\,/\,$0.7449$ on CR2, exceeding the compact PlainUNet
baseline by $+0.502$\,dB on CR1 and $+0.157$\,dB on CR2 despite
operating with sparse spike activations and TDBN-quantized
intermediate features. The smaller margin on CR2 reflects the harder
restoration ceiling under thick-cloud occlusion, where the encoder
features themselves carry less recoverable information.

\begin{table}[!htbp]
\centering
\caption{Centralized cloud-removal performance on CUHK-CR1 and
CUHK-CR2. \OrbitALIF{} is trained for $600$ epochs; PlainUNet is the
compact ANN U-Net baseline trained for $400$ epochs. PSNR and SSIM
are evaluated on $64\!\times\!64$ centre patches of the held-out
test split. Bold indicates the best per dataset.}
\label{tab:centralized}
\begin{tabular}{lcccc}
\toprule
Method & Params (M) & PSNR (dB) & SSIM & Wall (h) \\
\midrule
\multicolumn{5}{l}{\emph{CUHK-CR1 (thin clouds)}} \\
PlainUNet                  & $1.31$ & $24.872$           & $0.7683$           & $0.60$ \\
\OrbitALIF{} (ours)        & $2.30$ & $\mathbf{25.374}$  & $\mathbf{0.7728}$  & $2.30$ \\
\midrule
\multicolumn{5}{l}{\emph{CUHK-CR2 (thick clouds)}} \\
PlainUNet                  & $1.31$ & $23.040$           & $0.7400$           & $0.60$ \\
\OrbitALIF{} (ours)        & $2.30$ & $\mathbf{23.197}$  & $\mathbf{0.7449}$  & $2.34$ \\
\bottomrule
\end{tabular}
\end{table}

\begin{figure}[!htbp]
  \centering
  \includegraphics[width=\columnwidth]{fig7_centralized_4panel.pdf}
  \caption{Centralized training dynamics on CUHK-CR1 (top row) and
  CUHK-CR2 (bottom row). \OrbitALIF{} (blue) reaches a higher PSNR
  plateau than the compact PlainUNet baseline (orange) on both
  datasets despite operating with sparse spike activations. Loss
  panels (a) and (c) use a logarithmic ordinate.}
  \label{fig:fig7-centralized-4panel}
\end{figure}

The centralized advantage is architectural rather than capacity-driven:
\OrbitALIF{} uses $2.30$\,M parameters versus $1.31$\,M for PlainUNet,
but the $+0.50$\,dB CR1 gain stems from the four \OrbitALIF{}
building blocks rather than from raw parameter count, as the
ablation study in Section~\ref{sec:exp-ablation} confirms by
selectively disabling them.

\textbf{Capacity-scaled variant.}\;\;For completeness, we additionally
instantiate \OrbitALIF{} at $C\!=\!56$ ($\approx\!14$\,M parameters)
on the same dataflow to study the quality--computation trade-off.
A detailed comparison against full-capacity restoration baselines
(SpA-GAN, MemoryNet, Restormer, DDFR) is deferred to the extended
version because the corresponding centralized runs are still in
progress at submission time and the focus of this paper is the
onboard deployment regime.

\subsection{Component Ablation}\label{sec:exp-ablation}

We disable one \OrbitALIF{} component at a time and retrain on
CUHK-CR1 for $300$ epochs. Ablation B1 surgically replaces every
\SHAM{} block with the identity mapping; B2 replaces every \DFRB{}
with a single-group residual block that retains only the temporal
LIF group (and the shortcut, TCS-Att, and \SHAM{} tail) so that the
PixelShuffle-LIF spatial high-frequency path and the cross-scale gate
$\sigma(g)$ are removed; B3 replaces \FQS{} with a binary spike whose
forward output is $\{0,1\}$ but whose surrogate gradient window
$[0,4]$ and $1/4$ scale are kept identical to \FQS{}, so that only
the forward quantization granularity changes. \AGFM{} is not
separately ablated because its sigmoid gate degenerates to plain
addition at $\boldsymbol{\gamma}\!=\!\tfrac{1}{2}$ and no surgery is
needed in code.

\begin{table}[!htbp]
\centering
\caption{Component ablation on CUHK-CR1 ($300$ epochs).
$\Delta$ values are relative to the full \OrbitALIF{} (A1). Row
\emph{C2} reproduces the PlainUNet baseline of
Table~\ref{tab:centralized} for direct reference.}
\label{tab:ablation}
\begin{tabular}{lccccc}
\toprule
Variant & Params (M) & PSNR & $\Delta$\,PSNR & SSIM & $\Delta$\,SSIM \\
\midrule
\textbf{Full (A1)}        & $2.30$ & $\mathbf{25.374}$ & ---         & $\mathbf{0.7728}$ & ---           \\
$-$\SHAM{} (B1)           & $2.22$ & $24.777$         & $-0.597$    & $0.7620$         & $-0.0108$    \\
$-$\DFRB{} dual (B2)      & $1.44$ & $24.858$         & $-0.516$    & $0.7633$         & $-0.0095$    \\
$-$\FQS{} (B3, binary)    & $2.30$ & $25.072$         & $-0.302$    & $0.7666$         & $-0.0062$    \\
\midrule
PlainUNet (C2, ref.)      & $1.31$ & $24.872$         & $-0.502$    & $0.7683$         & $-0.0045$    \\
\bottomrule
\end{tabular}
\end{table}

\begin{figure}[!htbp]
  \centering
  \includegraphics[width=\columnwidth]{fig4_ablation_bars.pdf}
  \caption{Component ablation on CUHK-CR1. PSNR (left) and SSIM
  (right) for the full \OrbitALIF{} (A1) and the three
  single-component ablations B1--B3. Numerical $\Delta$ values are
  printed above each ablated bar; the diagonal-hatch pattern marks
  the reference (Full) bar.}
  \label{fig:ablation}
\end{figure}

\textbf{Component ranking.}\;\;Table~\ref{tab:ablation} attributes the
largest single-component PSNR drop to \SHAM{} ($-0.597$\,dB),
followed by the dual-group \DFRB{} ($-0.516$\,dB) and \FQS{}
($-0.302$\,dB). The ranking is consistent with the design rationale
of Section~\ref{sec:framework-arch}. \SHAM{} contributes orthogonal
recalibration along the temporal-amplitude and frequency-spatial
axes that no other component recovers. The dual-group \DFRB{}
provides the spatial high-frequency path that complements the
temporal LIF group; without it the model loses fine-texture
restoration capacity. \FQS{} contributes the smallest drop because
under the $300$-epoch budget the membrane potential takes
$\sim\!50$ epochs to warm up past $V_{\rm th}$ under TDBN
initialisation, leaving the multi-level dynamic range partially
under-utilised; under the full $600$-epoch schedule of the A1 run
this gap is expected to widen.

\textbf{Co-design evidence.}\;\;A striking observation is that two of
the three single-component ablations (B1 at $24.777$\,dB and B2 at
$24.858$\,dB) fall \emph{below} the PlainUNet baseline at
$24.872$\,dB, even though the ablated networks retain the rest of the
\OrbitALIF{} architecture. Removing one component therefore breaks
\OrbitALIF{} more than reverting entirely to a plain ANN. This is
strong evidence that the \OrbitALIF{} blocks are co-designed: the
$+0.50$\,dB CR1 advantage is not the sum of independently-additive
component gains but a synergistic effect that requires all
components to operate together. Only B3, which keeps the
architecture intact and only relaxes the spike granularity, remains
above PlainUNet, confirming that \FQS{} is a refinement rather than
a load-bearing component.

\subsection{Decentralized Federated Cloud-Removal}\label{sec:exp-federated}

This subsection evaluates the central system contribution of this
work: the decentralized \textsc{FedBN}\,+\,Gossip training pipeline of
Sections~\ref{sec:sysmodel-procedure} and~\ref{sec:framework-fl}. We
deploy three matched-protocol backbones on the
$5\!\times\!10\!=\!50$-satellite Walker-Star constellation under
source-level Dirichlet $\alpha\!=\!0.1$ for $R\!=\!200$ rounds.
\OrbitALIF{}-SNN is the full spiking variant; \OrbitALIF{}-ANN replaces
every \FQS{}-LIF with ReLU; PlainUNet is the compact ANN U-Net
baseline. All three share the identical aggregation pipeline of
Section~\ref{sec:sysmodel-procedure} and differ only in the
per-satellite network, so any performance gap is attributable to the
backbone alone.

\begin{table}[!htbp]
\centering
\caption{Decentralized federated cloud-removal on CUHK-CR1\,$\cup$\,CR2
under \textsc{FedBN}\,+\,Gossip, $R\!=\!200$ rounds, source-level
Dirichlet $\alpha\!=\!0.1$. Communication volume is the cumulative
inter-plane traffic over shared (non-BN) parameters. Bracketed cells
[\,$\cdot$\,] are pending the in-progress \OrbitALIF{}-SNN runs and
will be filled in the camera-ready or extended version.}
\label{tab:federated}
\begin{tabular}{lcccc}
\toprule
Backbone (FedBN $\times$ Gossip)    & PSNR (dB) & SSIM & Comm (MB) & Wall (h) \\
\midrule
PlainUNet (baseline)                 & $22.012$  & $0.6819$  & $8420$     & $4.37$ \\
\OrbitALIF{}-ANN                     & [\,XX.XX] & [\,0.XXXX]& [\,XX]     & [\,X.XX] \\
\OrbitALIF{}-SNN (ours)              & [\,XX.XX] & [\,0.XXXX]& [\,XX]     & [\,X.XX] \\
\bottomrule
\end{tabular}
\end{table}

\begin{figure}[!htbp]
  \centering
  \includegraphics[width=\columnwidth]{fig5_federated_curves.pdf}
  \caption{Federated test PSNR versus communication round under
  \textsc{FedBN}\,+\,Gossip with $\alpha\!=\!0.1$. Per-plane shading
  shows the min--max envelope across the five orbital planes. The
  spiking-backbone curve is deferred to the extended version.}
  \label{fig:federated}
\end{figure}

\textbf{Federated reference.}\;\;The PlainUNet baseline reaches
$22.012$\,dB\,/\,$0.6819$ at round $200$ with $8.42$\,GB cumulative
inter-plane traffic in $4.37$\,h on a single GPU. This constitutes a
\emph{federated PSNR floor} against which the \OrbitALIF{} variants
will be measured. The $2.86$\,dB gap to the centralized PlainUNet
result on CR1 ($24.872$\,dB) quantifies the cost of the source-level
Dirichlet $\alpha\!=\!0.1$ partition: under this concentration $30$
of $50$ satellites hold only the floor of five training tiles each,
so individual satellites cannot recover the full encoder
representation, and per-round Gossip mixing only partially
compensates over $R\!=\!200$ rounds.

\textbf{Dirichlet partition statistics.}\;\;The actual per-client
sample-count vector recorded by the partition log includes
$30/50\!=\!60\,\%$ floor clients (five tiles each) and three heavy
clients with $\{117, 118, 118\}$ samples; the remaining $17$ clients
are spread between these two extremes. The same $\alpha\!=\!0.1$
configuration is used for all federated runs in
Table~\ref{tab:federated}, so any backbone-level differences cannot
be attributed to differences in the data partition.

\subsection{Energy Analysis}\label{sec:exp-energy}

\begin{figure}[!htbp]
  \centering
  \includegraphics[width=\columnwidth]{fig6_energy_bars.pdf}
  \caption{Per-image estimated inference energy on a
  $64\!\times\!64$ CUHK-CR1 centre patch. ANN baseline at
  $4.6$\,pJ/MAC (Horowitz, ISSCC 2014, $45$\,nm CMOS); \OrbitALIF{}
  is reported as a dual bound: a \emph{conservative upper bound}
  that charges every non-zero MAC at the full ANN rate, and a
  \emph{deployment-target lower bound} at $77$\,fJ/SOP for
  neuromorphic substrates. Multipliers vs.\ ANN are annotated above
  each bar.}
  \label{fig:energy}
\end{figure}

\begin{figure}[!htbp]
  \centering
  \includegraphics[width=\columnwidth]{fig9_per_layer_spike_rate.pdf}
  \caption{\OrbitALIF{} per-layer LIF firing rate on CUHK-CR1 test,
  measured through forward hooks on every LIF / \FQS{} module.
  Binary LIF layers exhibit $20$\,--\,$70\,\%$ firing; \FQS{}-equipped
  layers show a higher non-zero floor because four of the five
  quantization levels are non-zero by construction.}
  \label{fig:firing-rate}
\end{figure}

\textbf{Methodology and disclosure.}\;\;Per-image inference energy is
\emph{estimated}, not hardware-measured. We hook every
\texttt{Conv2d}/\texttt{Linear} layer to count MACs and the input
non-zero rate $r_\ell$, and every LIF / \FQS{} module to record the
output firing rate, on $32$ CUHK-CR1 test patches at
$64\!\times\!64$. The MAC counts and per-layer rates are real
measurements; the per-operation energy constants are taken from the
literature. We use $4.6$\,pJ/MAC for the ANN baseline (Horowitz, ISSCC
2014, $45$\,nm CMOS) and $77$\,fJ/SOP for the SNN deployment estimate,
the constant adopted by FLSNN~\cite{ref:flsnn} and
ESDNet~\cite{ref:esdnet} for neuromorphic substrates such as
Loihi-2~\cite{ref:loihi}. Hardware-measured power on actual
neuromorphic platforms is deferred to deployment validation.

\textbf{Total energy bounds.}\;\;Figure~\ref{fig:energy} reports
per-image bounds. The ANN baseline at $4.6$\,pJ/MAC over the
$4.51$\,GMAC inference consumes $20.75$\,mJ. \OrbitALIF{} is reported
as a dual bound. The \emph{conservative upper bound} of $17.14$\,mJ
($1.21\!\times\!$ reduction) charges every non-zero MAC at the full
$4.6$\,pJ ANN rate and is achievable without any neuromorphic
hardware. The \emph{deployment-target lower bound} of
$0.287$\,mJ ($72.3\!\times\!$ reduction) uses $77$\,fJ/SOP and
applies to substrates that support the full SOP / spike abstraction
(e.g., Loihi-2, Speck, Akida). Across CR1 and CR2 inferences on the
same architecture, the measured effective non-zero MAC ratios are
$\bar{r}_{\rm CR1}\!=\!0.826$ and $\bar{r}_{\rm CR2}\!=\!0.823$,
yielding indistinguishable energy bounds at two decimal places. We
do not claim data-independence of the energy estimate; the two
CUHK-CR distributions are simply close enough to give consistent
bounds.

\textbf{Per-layer firing rates.}\;\;Figure~\ref{fig:firing-rate}
visualizes the layer-wise LIF / \FQS{} firing rates on CR1 test
patches. Binary LIF layers exhibit $20$\,--\,$70\,\%$ firing across
the analysis path. \FQS{}-equipped layers (notably the
\texttt{lif\_node} inside \SSHB{}) show a higher non-zero floor,
because four of the five \FQS{} levels $\{1/4,1/2,3/4,1\}$ are
non-zero by construction. The firing-rate measurement is the
empirical witness that spike sparsity is real rather than a
simulation artefact, and the per-layer non-zero MAC rates that enter
the energy formula
$E_{\rm SNN}\!=\!\sum_\ell\mathrm{MAC}_\ell\,r_\ell\,e_{\rm op}$
are derived from these measurements.

\subsection{Qualitative Results}\label{sec:exp-qualitative}

\begin{figure}[!htbp]
  \centering
  \includegraphics[width=\columnwidth]{fig3_qualitative_grid.pdf}
  \caption{Qualitative cloud-removal comparison on four representative
  CUHK-CR1 test tiles. Columns left-to-right: cloudy input,
  \OrbitALIF{} (ours), PlainUNet baseline, ground truth. Per-image
  PSNR is annotated on each restoration cell.}
  \label{fig:qualitative}
\end{figure}

Figure~\ref{fig:qualitative} compares centre-patch restorations on
four CUHK-CR1 test tiles selected via a deterministic
cloud-coverage severity score (the L2 distance between cloudy and
ground-truth pixels), biasing the selection toward harder samples.
Across all four tiles \OrbitALIF{} recovers texture and colour
consistency that visibly exceeds PlainUNet, particularly along
thin-cloud edges where the high-frequency content carried by the
\DFRB{} spatial path (Section~\ref{sec:framework-dfrb}) is most
active. The quality gap is concentrated in regions of moderate cloud
density (\AGFM{} gate $\boldsymbol{\gamma}\!\approx\!0.5$); under
fully-opaque cloud cover both architectures degrade similarly,
consistent with the smaller CR2 gain reported in
Table~\ref{tab:centralized}.

% =====================================================================
% Section V: Conclusion
% =====================================================================
\section{Conclusion}\label{sec:conclusion}

We presented \OrbitALIF{}, the first end-to-end decentralized
federated learning framework for satellite-onboard SNN-based optical
cloud removal. The framework couples an attention-enhanced LIF
restoration backbone with a constellation-native training pipeline.
The backbone integrates five-level quantized spikes (\FQS{}),
dual-frequency residual blocks (\DFRB{}), spectral-spatial hybrid
attention (\SHAM{}), adaptive gated cross-stage fusion (\AGFM{}),
and full-resolution spectro-temporal hyper-blocks (\SSHB{}) into a
sparse spiking dataflow that preserves pixel-level restoration
capability. The training pipeline runs entirely between satellites
through intra-plane ring AllReduce and inter-plane Gossip mixing,
with all BN-tagged parameters kept local under the \textsc{FedBN}
protocol; this avoids both the bandwidth cost of downlinking raw
observations and the dependence on a centralized ground parameter
server.

Empirically, the $2.30$\,M onboard deployment variant attains
$25.37$\,dB\,/\,$0.7728$ PSNR\,/\,SSIM on CUHK-CR1 and
$23.20$\,dB\,/\,$0.7449$ on CUHK-CR2, exceeding a compact PlainUNet
baseline by $+0.50$\,dB on CR1 and $+0.16$\,dB on CR2. Per-image
inference energy on a $64\!\times\!64$ patch is estimated as
$1.21\!\times$ to $72.3\!\times$ lower than the ANN baseline of
$20.75$\,mJ under standard literature pJ-per-operation constants.
Component ablations show that the building blocks form a co-designed
whole: removing \SHAM{} or the dual-group \DFRB{} drops PSNR
\emph{below} the plain ANN baseline, while removing \FQS{} (binary
spikes) keeps performance above PlainUNet. The centralized
PlainUNet-vs-federated PlainUNet gap of $2.86$\,dB on CR1 quantifies
the difficulty of the source-level Dirichlet $\alpha\!=\!0.1$
partition and frames the federated setting in which the spiking
backbone is evaluated.

Several directions extend this work. (i) Closing the
federated-vs-centralized gap, particularly under the spiking
backbone whose $\alpha$-sensitivity runs are deferred to the
extended version, by combining \textsc{FedBN}\,+\,Gossip with partial
participation, momentum aggregation, or learnable inter-plane
mixing matrices that exploit orbital geometry. (ii) Real-orbit
validation on neuromorphic substrates (Loihi-2, Speck, Akida) to
replace the analytic energy bracket with measured power draw under
realistic on-orbit thermal and power conditions. (iii) Multi-source
generalization beyond CUHK-CR1\,+\,CR2 (e.g., Sentinel-2 versus
Landsat-8) with $S\!\geq\!3$ source domains, exercising stronger
source-level heterogeneity. (iv) Capacity scaling of \OrbitALIF{} to
the $14$\,M variant and a full comparison against high-capacity
restoration baselines, characterising the quality--computation
trade-off across the full deployment spectrum.

% =====================================================================
% Bibliography
% =====================================================================
% IMPORTANT: Each \bibitem is tagged with one of three statuses:
%   [VERIFIED]    — bibliographic fields are stable across multiple
%                   independent sources (canonical / well-cited paper).
%                   Safe to keep as-is.
%   [USER VERIFY] — the paper exists in the public record but specific
%                   fields (page numbers, volume, doi, exact title
%                   capitalisation) should be cross-checked against the
%                   actual published / arXiv PDF before submission.
%   [PLACEHOLDER] — bibliographic skeleton only; user MUST replace with
%                   a real, verified citation.  Do not submit while any
%                   [PLACEHOLDER] entry remains.
% =====================================================================

\begin{thebibliography}{99}

% [PLACEHOLDER]
% Citation context: opening sentence, "real-time applications such as
% natural disaster monitoring, emergency response, agricultural
% observation...".  Recommendation: replace with a recent IEEE
% Geoscience and Remote Sensing Magazine survey on optical remote-
% sensing applications.  Sample candidates the user may consider
% (must be cross-checked individually):
%   - "Deep learning in remote sensing: A comprehensive review and
%      list of resources", Zhu et al., IEEE GRSM 2017.
%   - "Deep Learning Methods for Optical Remote Sensing Image
%      Restoration: A Survey", Liu et al., IEEE GRSM (year varies).
\bibitem{ref:earthobs1}
[PLACEHOLDER -- replace with verified survey on remote-sensing
applications.]

% [PLACEHOLDER]
% Citation context: same opening sentence, second authority citation.
% Recommendation: a survey on real-time satellite imagery for disaster
% / emergency response (e.g., IEEE GRSM or ISPRS Journal).
\bibitem{ref:earthobs2}
[PLACEHOLDER -- replace with verified survey on real-time satellite
imagery for disaster / emergency response.]

% [USER VERIFY]
% Title and venue match the public CVPRW 2024 (AI4Space workshop)
% paper "Tackling the Satellite Downlink Bottleneck with Federated
% Onboard Learning of Image Compression" by Gomez et al.  The CVF
% open-access page is at openaccess.thecvf.com (CVPRW 2024 AI4Space).
% USER ACTION: open the CVF PDF and confirm authors / page range /
% exact title before submission.
\bibitem{ref:onboardfl-bottleneck}
P.~Gomez \emph{et~al.},
``Tackling the satellite downlink bottleneck with federated onboard
learning of image compression,''
\emph{Proc.\ IEEE/CVF Conf.\ Comput.\ Vis.\ Pattern Recognit.\
Workshops (CVPRW)}, 2024.

% [PLACEHOLDER]
% Citation context: Space-CPN concept and architecture.  Recent
% Chinese-language and IEEE Network articles discuss this concept.
% USER ACTION: replace with one or two specific Space-CPN architecture
% papers (IEEE Network / IEEE Comm.\ Magazine).  Make sure the cited
% paper actually defines or surveys Space Computing Power Networks.
\bibitem{ref:spacecpn1}
[PLACEHOLDER -- replace with verified Space Computing Power Networks
architecture / survey paper.]

% [USER VERIFY]
% This is the CUHK-CR / DDFR paper.  An arXiv version is available at
% arXiv:2401.15105 with title
%   "Diffusion Enhancement for Cloud Removal in Ultra-Resolution
%    Remote Sensing Imagery"
% and the published version is in IEEE TGRS.  USER ACTION: confirm
% the final IEEE TGRS volume / issue / pages and the canonical author
% list against the arXiv abstract page or IEEE Xplore record.
\bibitem{ref:cuhkcr}
Y.~Sui \emph{et~al.},
``Diffusion enhancement for cloud removal in ultra-resolution remote
sensing imagery,''
\emph{IEEE Trans.\ Geosci.\ Remote Sens.}, 2024.
% arXiv:2401.15105

% [VERIFIED]
% Canonical SNN reference establishing spiking neurons as the third
% generation of neural network models.  Bibliographic fields are
% stable across decades of citations.
\bibitem{ref:maass1997}
W.~Maass,
``Networks of spiking neurons: The third generation of neural network
models,''
\emph{Neural Networks}, vol.~10, no.~9, pp.~1659--1671, 1997.

% [USER VERIFY]
% Loihi neuromorphic processor.  USER ACTION: confirm exact
% bibliographic fields (volume / issue / month) of IEEE Micro.
\bibitem{ref:loihi}
M.~Davies \emph{et~al.},
``Loihi: A neuromorphic manycore processor with on-chip learning,''
\emph{IEEE Micro}, vol.~38, no.~1, pp.~82--99, Jan./Feb.\ 2018.

% [USER VERIFY]
% ESDNet for SNN-based image deraining; arXiv preprint listed at
% arXiv:2207.02094.  USER ACTION: confirm whether a peer-reviewed
% version exists in a specific venue and update accordingly.
\bibitem{ref:esdnet}
H.~Song \emph{et~al.},
``Exploring the potentials of spiking neural networks for image
deraining,''
\emph{arXiv preprint arXiv:2207.02094}, 2022.

% [USER VERIFY]
% FLSNN: decentralized federated SNN over Walker-Star constellation
% for EuroSAT classification.  arXiv preprint identifier listed in
% the user's existing draft as arXiv:2501.15995, 2025.  USER ACTION:
% open the arXiv abstract page, copy the exact title and author list,
% and confirm the preprint ID.  This is a load-bearing citation for
% the related-work paragraph; treat it as critical.
\bibitem{ref:flsnn}
P.~Yang, T.~Wang, \emph{et~al.},
``Brain-inspired decentralized satellite learning in space computing
power networks,''
\emph{arXiv preprint arXiv:2501.15995}, 2025.

% [VERIFIED]
% FedAvg.  Canonical federated learning baseline; bibliographic fields
% are widely consistent.
\bibitem{ref:fedavg}
H.~B.~McMahan, E.~Moore, D.~Ramage, S.~Hampson, and B.~A.~y~Arcas,
``Communication-efficient learning of deep networks from decentralized
data,''
\emph{Proc.\ Int.\ Conf.\ Artif.\ Intell.\ Statist.\ (AISTATS)}, 2017.

% [VERIFIED]
% FedBN.  ICLR 2021.  Title and authorship are widely cross-cited.
\bibitem{ref:fedbn}
X.~Li, M.~Jiang, X.~Zhang, M.~Kamp, and Q.~Dou,
``FedBN: Federated learning on non-IID features via local batch
normalization,''
\emph{Proc.\ Int.\ Conf.\ Learn.\ Represent.\ (ICLR)}, 2021.

% [VERIFIED]
% Boyd et al. on randomized gossip algorithms.  Canonical reference
% for gossip-style decentralized averaging.
\bibitem{ref:gossip-boyd}
S.~Boyd, A.~Ghosh, B.~Prabhakar, and D.~Shah,
``Randomized gossip algorithms,''
\emph{IEEE Trans.\ Inf.\ Theory}, vol.~52, no.~6, pp.~2508--2530,
Jun.\ 2006.

\end{thebibliography}

\end{document}
